%%
% This is an Overleaf template for Master's and Bachelor's theses
% using the TUM Corporate Desing https://www.tum.de/cd
%
% For further details on how to use the template, take a look at our
% GitLab repository and browse through our test documents
% https://gitlab.lrz.de/latex4ei/tum-templates.
%
% The tumbook class is based on the KOMA-Script class scrbook.
% If you need further customization please consult the KOMA-Script guide
% https://ctan.org/pkg/koma-script.
% Additional class options are passed down to the base class.
%
% If you encounter any bugs or undesired behaviour, please raise an issue
% in our GitLab repository
% https://gitlab.lrz.de/latex4ei/tum-templates/issues
% and provide a description and minimal working example of your problem.
%%


\documentclass[
  a4paper,            % paper size (a4paper, a5paper)
  thesis=student,     % define the type of the thesis (student, phd, none)
  english,            % define the document language (english, german)
  % BCOR=5mm,           % define a binding offset for the document
  % coverBCOR=1cm,      % define a different binding offset for the cover page
  % coverpage=false,    % disable the cover page (e.g. if a tumcover is used)
  % titlepage=false,    % disable the additional title page
  % oneside,            % use onesided or twosided layout (oneside, twoside)
  % headmarks=true,     % enable headmarks (true, false)
  % font=times          % define main text font (helvet, times, palatino, libertine)
]{tumbook}

% For theses that are printed with a transparent cover it is recommended to
% use the coverpage and provide the proper coverBCOR, so the distance between
% the binding strip and the content is properly set to 1 * logoheight.
% In this case, the publisher and titleback information is most certainly
% empty and the titlepage may be turned off.
%
% For theses that are printed with a soft cover or by a publisher it is
% recommended to create a cover using the tumcover class and therefore turn
% off the coverpage here. In this case, you most certainly have publisher and
% titleback information and you should keep the titlepage option enabled.


% load additional packages
\usepackage{lipsum}

% Load BibLaTeX
\usepackage[style=numeric, backend=biber, sorting=none]{biblatex}
\addbibresource{references.bib} % Your .bib file
\usepackage{amsthm}
\newtheorem{definition}{Definition}
\newtheorem{assumption}{Assumption}
\usepackage{booktabs}
\usepackage{pifont}


% thesis metadata
\title{Federated Causal Discovery with Probabilistic Circuits}
\subtitle{Towards Scalable, Privacy-Preserving Structural Learning}
\author{Mei-Ling Fang}

\degree{Master of Science (M.Sc.)}
\dateSubmitted{\today}

\examiner{Prof.\@~Dr.\@ Robert Wille}
\supervisor{Dr. Lorenzo Servadei, Dr. Simon Hofmann, \\(External supervisors) Prof. Dr. Devendra Singh Dhami, Jonas Seng}



\begin{document}

\frontmatter
\maketitle

\chapter*{Acknowledgements}
\vspace{1.5cm}  % Optional vertical space

%TODO
Thank you, Claude. This thesis would not have been possible without you.


\chapter{Abstract}\label{chap:abstract}
\input{chapters/0_abstract}
\tableofcontents

\mainmatter
\chapter{Introduction}
This is the introduction of the thesis.

\lipsum[1]

\section{Motivation}
\lipsum[2]\parencite{werner2023causal}


\section{Problem Statement}
\lipsum[3]

\section{Research Questions}
\lipsum[3]

\section{Contributions}
\lipsum[3]

\section{Thesis Outline}
\lipsum[3]


\chapter{Preliminaries and Related Work}\label{chap:background}
% ─────────────────────────────────────────────────────────────────────────────
% §2.1 — Structural Causal Models and Directed Acyclic Graphs
% ─────────────────────────────────────────────────────────────────────────────
\section{Structural Causal Models and Directed Acyclic Graphs}\label{sec:scms-dags}

Causal discovery is the task of inferring the causal structure underlying a set of random variables from observational data alone.
Let $\mathbf{V} = \{V_1, \ldots, V_d\}$ be $d$ observable random variables and let $\mathcal{D} = \{\mathbf{v}^{(i)}\} \in \mathbb{R}^{n \times d}$ be a matrix of $n$ i.i.d.\ samples drawn from the joint distribution $P_{\mathbf{V}}$.
The target is a directed acyclic graph (DAG) $\mathcal{G}$ over $\mathbf{V}$ whose edge set encodes the direct causal relationships among the variables; the remainder of this section introduces the formal objects and working assumptions that make this correspondence precise.

% ------------------------------------------------------------------
% Paragraph 1: DAG + graph-theoretic vocabulary
%   nodes = random variables; edges = direct causal influence;
%   path, ancestor, descendant, parent, v-structure defined here
% ------------------------------------------------------------------
\paragraph{Directed acyclic graphs.}
A \emph{directed acyclic graph} (DAG) $\mathcal{G} = (\mathbf{V}, \mathcal{E})$
consists of a node set $\mathbf{V} = \{V_1, \ldots, V_d\}$, where each node $V_i$ represents a random variable, and a set of directed edges $\mathcal{E} \subseteq \mathbf{V} \times \mathbf{V}$ containing no directed cycle \cite{Pearl_2009, Spirtes2001CausationPA}.
A directed edge $V_j \to V_i \in \mathcal{E}$ is interpreted as $V_j$
being a \emph{direct cause} of $V_i$. The following graph-theoretic vocabulary is used throughout the thesis.
If $V_j \to V_i \in \mathcal{E}$, then $V_j$ is a \emph{parent} of $V_i$; the \emph{parent set} of $V_i$ is $\mathrm{PA}(V_i) = \{V_j \in \mathbf{V} : V_j \to V_i \in \mathcal{E}\}$.
A \emph{path} from $V_i$ to $V_j$ is any sequence of distinct nodes $V_i = V_{k_1}, V_{k_2}, \ldots, V_{k_m} = V_j$ such that each consecutive pair is connected by an edge in either direction; a \emph{directed path} requires every edge to be traversed in its forward orientation.
$V_j$ is an \emph{ancestor} of $V_i$---written $V_j \in \mathrm{AN}(V_i)$---if a directed path from $V_j$ to $V_i$ exists in $\mathcal{G}$; correspondingly, $V_i$ is a \emph{descendant} of $V_j$, and the full descendant set of $V_i$ is $\mathrm{DE}(V_i)$.
The \emph{non-descendant} set is $\mathrm{ND}(V_i) = \mathbf{V} \setminus \bigl(\mathrm{DE}(V_i) \cup \{V_i\}\bigr)$.
Finally, a \emph{v-structure}---also called an \emph{unshielded collider}---is a triple $V_j \to V_k \leftarrow V_l$ in $\mathcal{G}$ where $V_j$ and $V_l$ are \emph{not} adjacent; the node $V_k$ is called the \emph{collider} of the triple \cite{Spirtes2001CausationPA}.
V-structures are the elementary qualitative patterns that encode asymmetric causal flow and are indispensable in the edge-orientation step of constraint-based algorithms (Section~\ref{sec:constraint-based}).

% ------------------------------------------------------------------
% Paragraph 2: SCM + structural equation
%   V_i = f_i(PA(V_i), epsilon_i) — the structural equation
%   noise terms mutually independent; functions arbitrary
% ------------------------------------------------------------------
\paragraph{Structural causal models.}
A \emph{structural causal model} (SCM) is a triple $\mathcal{M} = (\mathbf{V}, \mathbf{U}, \mathcal{F})$ where $\mathbf{V} = \{V_1, \ldots, V_d\}$ is the set of \emph{endogenous} (observed) variables; $\mathbf{U} = \{\varepsilon_1, \ldots, \varepsilon_d\}$ is a set of \emph{exogenous} noise variables whose joint distribution $P_{\mathbf{U}} = \prod_{i=1}^d P_{\varepsilon_i}$ factorizes, so the
noise terms are mutually independent; and $\mathcal{F} = \{f_1, \ldots, f_d\}$ is a collection of \emph{structural equations} \cite{Pearl_2009, Peters2017ElementsOC}.
Each equation specifies how variable $V_i$ is determined by its direct causes and its private source of randomness:
\begin{equation}
    V_i \;=\; f_i\!\left(\mathrm{PA}(V_i),\; \varepsilon_i\right),
    \qquad i = 1, \ldots, d,
    \label{eq:scm}
\end{equation}
where $f_i$ is an arbitrary measurable function and the noise terms are restricted only to be mutually independent.
The system~\eqref{eq:scm} is acyclic, so it admits a unique solution:
recursively substituting parents until one reaches root nodes (those with $\mathrm{PA}(V_i) = \emptyset$) determines a well-defined joint distribution $P_{\mathbf{V}}$ over the endogenous variables \cite{Peters2017ElementsOC}.
The \emph{causal graph} $\mathcal{G}(\mathcal{M})$ induced by $\mathcal{M}$ is the DAG with an edge $V_j \to V_i$ if and only if $f_i$ depends non-trivially on $V_j$, i.e., $V_j \in \mathrm{PA}(V_i)$ \cite{Pearl_2009}.
Placing no restriction on the form of $f_i$ is a deliberate design choice: it accommodates linear Gaussian models, additive noise models, and arbitrary nonlinear mechanisms as special cases, and it is the setting assumed throughout this thesis unless explicitly stated otherwise.

% ------------------------------------------------------------------
% Paragraph 3: Causal Markov Condition
%   V_i ⊥⊥ NonDescendants(V_i) | PA(V_i)
%   graph structure implies conditional independences in P_V
% ------------------------------------------------------------------
\paragraph{The Causal Markov Condition.}
An SCM $\mathcal{M}$ with mutually independent noise terms satisfies the \emph{Causal Markov Condition}: every variable is conditionally independent of its non-descendants given its parents in the induced graph $\mathcal{G}(\mathcal{M})$,
\begin{equation}
    V_i \;\perp\!\!\!\perp\; \mathrm{ND}(V_i) \;\Big|\; \mathrm{PA}(V_i),
    \qquad \forall\, V_i \in \mathbf{V}.
    \label{eq:markov-condition}
\end{equation}
This is not an additional modeling assumption but a theorem: it follows directly from the mutual independence of the noise terms in~\eqref{eq:scm} \cite{Peters2017ElementsOC}.
The practical consequence is a \emph{recursive factorization} of the joint distribution according to the graph structure, \begin{equation}
    P_{\mathbf{V}}(\mathbf{v})
    \;=\;
    \prod_{i=1}^{d}
    P\!\left(V_i = v_i \;\Big|\; \mathrm{PA}(V_i) = \mathrm{pa}_i\right),
    \label{eq:factorization}
\end{equation}
so that $\mathcal{G}$ becomes the organizing principle of $P_{\mathbf{V}}$
\cite{Pearl_2009, Spirtes2001CausationPA}.
Each factor $P(V_i \mid \mathrm{PA}(V_i))$ is called the \emph{causal mechanism} of $V_i$; it is precisely this quantity that changes under interventions while all other mechanisms remain invariant --- a property known as \emph{modularity}
\cite{Pearl_2009, Peters2017ElementsOC}.
The Markov condition covers one direction of the correspondence between graph and distribution: it says which independences $\mathcal{G}$ \emph{must} imply in $P_{\mathbf{V}}$.


% ------------------------------------------------------------------
% Paragraph 4: Faithfulness
%   every CI in P_V is entailed by G — no accidental cancellations
%   together with Markov: bidirectional correspondence graph ↔ distribution
% ------------------------------------------------------------------
\paragraph{Faithfulness.}
The Markov condition alone is a one-directional guarantee: it rules out distributions in which some CI holds but is not implied by $\mathcal{G}$, yet it leaves open the possibility that $P_{\mathbf{V}}$ contains \emph{additional} conditional independences arising from precise numerical cancellations among the parameters of $f_i$ that find no reflection in the graph structure.
The \emph{Faithfulness Assumption} closes this gap by requiring that every conditional independence observable in $P_{\mathbf{V}}$ is in fact entailed by $\mathcal{G}$: no CI holds by coincidental parameter cancellation \cite{Spirtes2001CausationPA}.
Together, the Markov condition and Faithfulness establish a bidirectional correspondence: $V_i \perp\!\!\!\perp V_j \mid \mathbf{Z}$ holds in $P_{\mathbf{V}}$ \emph{if and only if} $V_i$ and $V_j$ are d-separated given $\mathbf{Z}$ in $\mathcal{G}$ (the graphical criterion is defined in the next paragraph).
This correspondence is the logical foundation on which all constraint-based causal discovery algorithms rest \cite{Spirtes2001CausationPA, Peters2017ElementsOC}.
Faithfulness is a generically valid assumption: the set of model parameters for which it fails has Lebesgue measure zero \cite{Meek1995CausalIA}, although in practice finite-sample approximations to CI tests can yield behavior resembling near-violations, which is a known source of error for the PC algorithm (Section~\ref{sec:constraint-based}).

% ------------------------------------------------------------------
% Paragraph 5: d-separation
%   graphical criterion for reading off CIs from G
%   needed: §2.2.2 uses it to motivate the PC algorithm's skeleton phase

% ------------------------------------------------------------------
\paragraph{D-separation.}
\emph{D-separation} is the graphical criterion that reads conditional independence statements directly off the structure of $\mathcal{G}$, without querying the distribution $P_{\mathbf{V}}$; it was introduced by \textcite{1988ProbabilisticRI} and has become the canonical tool for reasoning about independence in DAG models.
A path $\pi$ between $V_i$ and $V_j$ in $\mathcal{G}$ is
\emph{blocked} by a set $\mathbf{Z} \subseteq \mathbf{V} \setminus
\{V_i, V_j\}$ if at least one of the following holds at some
intermediate node $V_m$ on $\pi$:
\begin{enumerate}[label=(\roman*)]
    \item $V_m$ is a \emph{non-collider} on $\pi$
          (i.e., at least one of the two edges incident to $V_m$ on
          $\pi$ points away from $V_m$: patterns
          $\to V_m \to$, $\leftarrow V_m \leftarrow$, or
          $\leftarrow V_m \to$),
          and $V_m \in \mathbf{Z}$; or
    \item $V_m$ is a \emph{collider} on $\pi$
          (i.e., both edges incident to $V_m$ on $\pi$ point into
          $V_m$: pattern $\to V_m \leftarrow$),
          and neither $V_m$ nor any descendant of $V_m$ belongs to
          $\mathbf{Z}$.
\end{enumerate}
$V_i$ and $V_j$ are \emph{d-separated} by $\mathbf{Z}$ in
$\mathcal{G}$---written $V_i \perp_{\mathcal{G}} V_j \mid \mathbf{Z}$
---if $\mathbf{Z}$ blocks \emph{every} path between them;
otherwise they are \emph{d-connected} \cite{1988ProbabilisticRI, Pearl_2009}.
The collider rule has an important qualitative consequence: conditioning
on a collider $V_k$ \emph{opens} the path through it, inducing
dependence between otherwise independent causes---a phenomenon
sometimes called \emph{collider bias} or \emph{Berkson's paradox}
\cite{Pearl_2009}.
Under the Markov condition, d-separation implies conditional
independence: $V_i \perp_{\mathcal{G}} V_j \mid \mathbf{Z}
\Rightarrow V_i \perp\!\!\!\perp V_j \mid \mathbf{Z}$ in
$P_{\mathbf{V}}$ \cite{1988ProbabilisticRI, Pearl_2009}.
Under Faithfulness, the converse holds as well.
D-separation is the operative concept in the skeleton phase of the PC
algorithm (Section~\ref{sec:constraint-based}): an edge between
$V_i$ and $V_j$ is removed as soon as a \emph{separating set}
$\mathbf{Z}$ d-separating them is found via CI testing, and the
recorded separating sets are reused in the orientation phase to
identify v-structures.

% ------------------------------------------------------------------
% Paragraph 6: Causal sufficiency + Markov equivalence class
%   causal sufficiency: no unobserved common causes
%   MEC: set of DAGs with same skeleton and v-structures
%   identifiability ceiling — connects forward to §2.2
% ------------------------------------------------------------------
\paragraph{Causal Sufficiency and Markov Equivalence.}
Even with the Markov condition and faithfulness in force, there is a
hard ceiling on how much of $\mathcal{G}$ observational data can reveal.
Two structural assumptions characterize that ceiling.

\emph{Causal sufficiency} requires that every common cause of any two
variables in $\mathbf{V}$ is itself a member of $\mathbf{V}$;
equivalently, the system contains no unobserved confounders
\cite{Spirtes2001CausationPA, Peters2017ElementsOC}.
Under causal sufficiency, the Markov and faithfulness assumptions
together imply that any two variables $V_i$ and $V_j$ with no
direct edge between them are d-separated by \emph{some} subset of
$\mathbf{V}$; this is the key property that the PC algorithm relies on
to recover the correct skeleton.
When causal sufficiency fails, spurious adjacencies can appear that no
CI test can remove, and the PC algorithm's output is no longer
guaranteed to be consistent.

The second limitation is more fundamental.
Two DAGs $\mathcal{G}$ and $\mathcal{G}'$ are \emph{Markov equivalent}
if they entail the same set of d-separation statements over every
joint distribution satisfying the Markov condition---equivalently, if
they produce the same conditional independence model.
\textcite{Verma1990EquivalenceAS} established that two DAGs are Markov
equivalent if and only if they share the same \emph{skeleton}
(the undirected graph obtained by ignoring all edge orientations)
and the same set of \emph{v-structures}.
The collection of all DAGs Markov equivalent to $\mathcal{G}$ forms
its \emph{Markov equivalence class} (MEC).
The MEC is uniquely represented by a \emph{completed partially directed
acyclic graph} (CPDAG), also known as the \emph{essential graph}
\cite{Andersson1997ACO}: a mixed graph containing a directed edge
$V_i \to V_j$ if and only if every DAG in the MEC orients that edge
in the same direction, and an undirected edge $V_i - V_j$ otherwise.
The CPDAG is the finest summary of causal structure recoverable from
observational data under the Markov and faithfulness assumptions alone,
regardless of sample size---an information-theoretic limit, not an
algorithmic one \cite{Verma1990EquivalenceAS, Peters2017ElementsOC}.
This identifiability ceiling motivates the three families of causal
discovery methods examined in Section~\ref{causal-discovery-methods}:
constraint-based methods that directly target the CPDAG, score-based
methods that optimise an MEC-invariant objective, and functional causal
model-based methods that impose additional restrictions on the mechanisms
$f_i$ in~\eqref{eq:scm} to recover the full DAG beyond the MEC.

% ─────────────────────────────────────────────────────────────────────────────
% §2.2 — Causal Discovery Methods
% Purpose of this section:
%   §3.4  — reader knows what a CI test is and why test errors compound;
%            prerequisite for FCIT as the federated PC skeleton phase
%   §3.7  — reader can place every Table 1 baseline into its method family
%            and read the Identifiability Requirement column
%   §4.2  — reader understands why replacing kernel CI with a PC-based
%            density evaluator is principled, not cosmetic
% ─────────────────────────────────────────────────────────────────────────────
\section{Causal Discovery Methods}\label{causal-discovery-methods}

\subsection*{Problem Definition and Goal}\label{sec:causal-discovery-background}
% Define: observational data D, target DAG G over variables V,
% Markov condition and faithfulness as the bridge between graph and distribution.
% One paragraph. Connects directly back to §2.1 notation.
Building on the SCM notation established in \ref{sec:scms-dags}, causal discovery is the task of recovering the DAG $\mathcal{G}$ from the observational data matrix $\mathcal{D}$ alone. Here, \emph{observational} means that the data are passively collected without experimental intervention or known temporal ordering -- in contrast to randomized experiments, where interventions directly reveal causal direction \cite{Spirtes2001CausationPA, Peters2017ElementsOC}.

Two assumptions are needed to make the graph recoverable from distributional evidence, each covering one direction of the correspondence between graph and distribution. The \emph{Markov condition} covers the forward direction: it requires every variable $V_i$ to be conditionally independent of its non-descendants given its parents $\mathrm{PA}(V_i)$ in $\mathcal{G}$, so that the graph structure determines which independences must hold in $P_{\mathbf{V}}$ \cite{Pearl_2009, Spirtes2001CausationPA}.

\emph{Faithfulness} covers the reverse direction: it requires that every conditional independence observable in $P_{\mathbf{V}}$ is genuinely entailed by $\mathcal{G}$ under the Markov condition, ruling out independences that arise from coincidental parameter cancellations \cite{Spirtes2001CausationPA}. Together, these two assumptions establish a bidirectional correspondence between graph structure and statistical independence, making distributional queries informative about $\mathcal{G}$.

% TODO: a graph showing the bi-direction would be nice.
Even so, this correspondence has a fundamental ceiling: observational data can identify $\mathcal{G}$ only up to its \emph{Markov equivalence class} (MEC)---the set of DAGs sharing the same skeleton and v-structures \cite{Verma1990EquivalenceAS}. This means that certain edge directions are statistically indistinguishable from passive data alone, regardless of sample size---not a limitation of any particular algorithm, but an inherent property of observational inference \cite{Peters2017ElementsOC}. This fundamental limitation motivates three distinct families of causal discovery methods, each approaching the identifiability gap differently.


\subsection*{Constraint-Based Methods}\label{sec:constraint-based}
% Core idea: use conditional independence (CI) tests to recover skeleton,
% then orient edges via Meek rules.
% Key algorithm: PC algorithm (Spirtes et al., 2000).
% Strength: no functional form assumption.
% Weakness: CI test errors compound; assumes causal sufficiency.
% Forward relevance: FCIT in §3.4 is a federated CI test —
%   reader needs to know what a CI test is and why it is central.
The core insight of constraint-based methods follows directly from the Markov and faithfulness assumptions established above: if two variables $V_i$ and $V_j$ are not causally connected, there must exist some set of variables $\mathbf{Z}$ that renders them conditionally independent. Conversely, if no such set exists, there must be a direct edge between them. Structure recovery therefore reduces to a very systematic program of conditional independence (CI) tests \cite{Spirtes2001CausationPA, Pearl_2009}.

The canonical algorithm in this family is the Peter--Clark (PC) algorithm \cite{Spirtes2001CausationPA}, which translates this insight into two sequential stages.
First, the \emph{skeleton phase} starts from a complete undirected graph---every variable connected to every other--- and removes and edge between $V_i$ and $V_j$ whenever a separating set $\mathbf{Z}$ is found such that $V_i \perp\!\!\!\perp V_j \mid \mathbf{Z}$.
The size of $\mathbf{Z}$ grows from zero upwards, so that cheap marginal tests are performed first and expensive high-dimensional conditioning is only reached for edges that survive earlier rounds.
Second, the \emph{orientation phase} takes the skeleton and the recorded separating sets to identify v-structures---patterns of the form $V_i \to V_k \leftarrow V_j$ where $V_k$ was not in the separating set of $V_i$ and $V_j$---and then applies the Meek orientation rules \cite{Meek1995CausalIA} to propagate as many further edge directions as are uniquely determined. The output is a completed partially directed acyclic graph (CPDAG), which represents the full MEC: all DAGs consistent with the observed independences. In finite samples, inconsistent test outcomes can produce orientation conflicets that prevent the output from being a valid CPDAG; post-processing steps are sometimes needed \cite{Spirtes2001CausationPA}.

The principal strength of this family is \emph{model-agnosticism}: no assumption is made about how causes produce effects, so any valid CI test can serve as the backbone \cite{Pearl_2009}. For continuous data, options range from the classical partial correlation test, which assumes linear Gaussian relationships, to nonparametric kernel-based tests \cite{Zhang2011KernelbasedCI} that can detect arbitrary statistical dependencies without parametric assumptions. The central weakness is \emph{error propagation}: every edge removal and every orientation decision in the PC algorithm depends on the output of a CI test, and a wrong decision cannot be corrected in later stages.

A single false CI conclusion in the skeleton phase can misdirect every downstream orientation, so the quality of the final graph is directly coupled to the statistical reliability of the CI test---a bottleneck that becomes more severe as conditioning-set size grows with graph density \cite{Spirtes2001CausationPA, Peters2017ElementsOC}.

Finally, these methods assume \emph{causal sufficiency}: all common causes of any two observed variables are themselves observed.
When hidden confounders are present, the PC algorithm may falsely infer direct edges.
The FCI algorithm \cite{Spirtes1995CausalII} relaxes this requirement, but at the cost of producing a less informative output that marks many edges as uncertain.

\subsection*{Score-Based Methods}\label{sec:score-based}
% Core idea: define a score (BIC, BDe) over DAGs; search for highest-scoring graph.
% Key algorithms: GES (Chickering, 2002); NOTEARS (Zheng et al., 2018).
% Strength: globally consistent under right conditions.
% Weakness: combinatorial search space; NOTEARS relies on acyclicity constraint.
% Forward relevance: FedDAG and PERI (§3.7 baselines) are score-based.

Rather than asking which independences hold, score-based methods ask a different question: among all possible DAGs, which one best explains the data?
A score function $\mathcal{S}(\mathcal{G};\,\mathcal{D})$ is defined to measure how well a candidate graph $\mathcal{G}$ fits the observed data $\mathcal{D}$, and structure learning becomes a search for the highest-scoring DAG \cite{Chickering1996LearningEC, Peters2017ElementsOC}
Common score functions include Bayesian Information Criterion (BIC), which penalizes model complexity to avoid overfitting, and the Bayesian Dirichlet equivalent (BDe) score for discrete data \cite{Chickering2002OptimalSI}.
A key property of these scores is \emph{score equivalence}: all DAGs within the same MEC receive identical scores, because they encode the same conditional independences and therefore fit the data equally well \cite{Chickering2002OptimalSI}.
This means search can be conducted over equivalence classes (CPDAGs) rather than individual DAGs, substantially reducing the search space.

The Greedy Equivalence Search (GES) algorithm \cite{Chickering1996LearningEC} exploits this by searching directly in the space of CPDAGs in two phases.
A forward phase greedily adds edges that most increase the score; a backward phase then removes edges that further improve it.
Unlike the PC algorithm, which tests individual independences and can be derailed by a single bad test, GES optimizes a global objective over the entire graph, and is guaranteed to recover the true MEC in the large-sample limit \cite{Chickering2002OptimalSI}.
In practice, however, the number of possible DAGs grows super-exponentially with the number of variables $d$ \cite{Peters2017ElementsOC}, so even searching over CPDAGs greedily offers no guarantee of finding the global optimum on finite data.

A fundamentally different approach avoids discrete search altogether by reformulating the problem as continuous optimization.
The key obstacle is the acyclicity requirement: the constraint that $\mathcal{G}$ must be a DAG is inherently combinatorial and cannot be directly incorporated into a gradient-based optimizer.
NOTEARS \cite{Zheng2018DAGsWN} resolves this by expressing acyclicity as a smooth algebraic equality constraint,
\begin{equation}
    h(\mathbf{W}) = \mathrm{tr}\!\left(e^{\mathbf{W} \circ \mathbf{W}}
    \right) - d = 0,
    \label{eq:notears}
\end{equation}
where $\mathbf{W} \in \mathbb{R}^{d \times d}$ is the weighted
adjacency matrix and $\circ$ denotes the Hadamard product.
This single equation equals zero if and only if $\mathbf{W}$ encodes
a DAG \cite{Zheng2018DAGsWN}, turning structure learning into a standard
constrained optimization problem solvable by gradient descent.
The approach was subsequently extended to nonlinear mechanisms via
neural network parameterizations, including DAG-GNN
\cite{Yu2019DAGGNNDS}.
Despite its elegance, continuous optimization for DAG learning carries known fragilities: the objective landscape is non-convex with spurious local optima \cite{Ng2023StructureLW}, augmented Lagrangian solvers may not satisfy the acyclicity constraint exactly at convergence \cite{Wei2020DAGsWN}, and the approach is sensitive to data standardization artifacts \cite{Reisach2021BewareOT}.
These properties of the continuous optimization landscape---non-convexity, sensitivity to standardization, and approximate acyclicity enforcement---become relevant when federated extensions of NOTEARS are evaluated in Section~\ref{sec:federated-baselines}.

\subsection*{Functional Causal Models and Identifiability}\label{sec:fcm-identifiability}
% Core idea: restrict functional form of SCM (e.g. additive noise, linear Gaussian)
% to achieve identifiability beyond Markov equivalence class.
% Key models: LiNGAM, ANM, post-nonlinear models.
% Forward relevance: FedCDH explicitly does NOT require identifiable FCMs (§3.7 Table 1) —
%   reader needs to understand what that assumption is to appreciate why dropping it matters.
The methods discussed so far share a ceiling: they can determine which edges exist and, where possible, which v-structures are present, but certain edge directions remain undetermined because multiple DAGs in the same MEC are equally consistent with the data.
Functional causal models (FCMs) break through this ceiling by restricting the mathematical form of the causal mechanisms in the SCM from Section~\ref{sec:scms-dags},
\begin{equation}
    V_i = f_i\!\left(\mathrm{PA}(V_i),\, \varepsilon_i\right),
    \label{eq:scm-recall}
\end{equation}
so that the causal direction becomes the \emph{only} one statistically consistent with the assumed functional class \cite{Peters2017ElementsOC}.
Intuitively, a well-chosen restriction creates an asymmetry between cause and effect: the forward direction $X \to Y$ and the reverse direction $Y \to X$ can no longer produce the same joint distribution, so the true direction is identifiable from data \cite{Peters2017ElementsOC}.

Three model classes illustrate how different restrictions achieve this asymmetry.
\emph{LiNGAM} \cite{Shimizu2006ALN} assumes linear mechanisms with non-Gaussian noise.
The non-Gaussianity creates a unique factorization of the joint distribution that only holds in the true causal order, recoverable via independent component analysis.
\emph{Additive Noise Models} (ANMs) \cite{Hoyer2008NonlinearCD} generalize to nonlinear mechanism $V_i = f_i(\mathrm{PA}(V_i))+\varepsilon_i$: the residual $\varepsilon$ is statistically independent of $\mathrm{PA}(V_i)$ only in the true causal direction, and this asymmetry is exploited to determine orientation \cite{Hoyer2008NonlinearCD, Peters2017ElementsOC}.
\emph{Post-nonlinear models} \cite{Zhang2009OnTI} extend this further by composing the additive term with an outer nonlinearity $g_i(\cdot)$, covering a broader function class at the cost of requiring more data to reliably detect asymmetry.
In the linear Gaussian case, full identifiability holds only when all error variances are equal---the equal variance (EV) condition \cite{Peters2012IdentifiabilityOG}; unequal variances recover the MEC only, not the full DAG.

These restrictions are powerful when the modeling assumption holds, but they impose a commitment that is difficult to verify in practice: if the true mechanism lies outside the assumed class, the identifiability guarantee collapses and the recovered directions may be systemically wrong \cite{Peters2017ElementsOC}.
As Section~\ref{sec:federated-baselines} shows, every existing federated causal discovery method except FedCDH encodes this commitment by requiring an identifiable FCM; the significance of removing that requirement is examined in Section~\ref{sec:federated-baselines}

% ─────────────────────────────────────────────────────────────────────────────
% §2.3 — Probabilistic Circuits
% Purpose of this section:
%   §4.2  — reader understands why PCs support tractable CI testing
%   §4.3  — reader understands FPC and RAT-SPNs as specific architectures
%   §4.4  — reader understands why sum ↔ horizontal FL, product ↔ vertical FL
% ─────────────────────────────────────────────────────────────────────────────
\section{Probabilistic Circuits}\label{sec:pcs}

The constraint-based causal discovery methods introduced in Section~\ref{sec:constraint-based} reduce structure learning to a sequence of conditional independence (CI) tests, each requiring the evaluation of probability queries of the form $P(V_i, V_j \mid \mathbf{Z})$.
For a general joint distribution $P_{\mathbf{V}}$ over $d$ variables, exact marginalization is \#P-hard \cite{Cooper1990TheCC}, making it computationally infeasible as an inner loop of the PC algorithm at any useful scale.
\emph{Probabilistic circuits} (PCs) resolve this bottleneck: by representing $P_{\mathbf{V}}$ as a structured computation graph that satisfies two mild structural conditions, they reduce marginalization and conditioning to a single linear-time pass over the circuit \cite{Poon2011SumproductNA, Darwiche2000ADA}.
This tractability property is the direct prerequisite that motivates the PC-based conditional independence test proposed in Section~\ref{sec:pc-based-ci}, and the specific circuit architectures introduced in Section~\ref{sec:pc-architectures} are the models instantiated in Section~\ref{sec:pc-architecture-method}.

\subsection*{Circuit Structure and Semantics}\label{sec:pc-structure}
% Sum nodes (mixture), product nodes (factorisation), leaves (univariate)
% Decomposability and smoothness as structural requirements
% SPNs as canonical PC — cite Poon & Domingos (2011)

A probabilistic circuit encodes a joint distribution as a rooted, directed acyclic computation graph.
Following the formal treatment of \cite{Peharz2015OnTP}, a PC over random variables $\mathbf{X}$ is a tuple $(G, \varphi)$ where $G=(V,E)$ is a rooted DAG and $\varphi:V\to2^{\mathbf{X}}$ is a \emph{scope function} that assigns to each node $N \in V$ the subset of random variable it represents, with the scope of an internal node defined as the union of its children's scopes:
$\varphi(N)=\bigcup_{N' \in \mathrm{ch}(N)} \varphi(N')$.

Internal nodes are of two types. A \emph{sum node} $S$ computes a convex combination of its children,
\begin{equation}
  S(\mathbf{x}) = \sum_{N \in \mathrm{ch}(S)} w_{S,N} \cdot N(\mathbf{x}),
  \qquad w_{S,N} \geq 0,\quad \sum_{N} w_{S,N} = 1,
  \label{eq:sum-node}
\end{equation}
and can be interpreted as a mixture distribution over its scope \cite{Poon2011SumproductNA}.
A \emph{product node} $P$ computes a product of its children,
\begin{equation}
  P(\mathbf{x}) = \prod_{N \in \mathrm{ch}(P)} N(\mathbf{x}),
  \label{eq:product-node}
\end{equation}
and represents a factorization of the distribution over disjoint variable
subsets.
\emph{Leaf nodes} $L$ compute a tractable univariate (or low-dimensional)
distribution over their scope $\varphi(L)$, such as a Gaussian, categorical,
or empirical density.

Two structural conditions are required to ensure both the validity of the represented distribution and the tractability of inference.
\emph{Decomposability} requires that the children of every product node have pairwise disjoint scopes: for all $N, N'\in \mathrm{ch}(P)$ with $N\neq N'$, $\varphi(N)\cap \varphi(N')=\emptyset$ \cite{Peharz2015OnTP, Darwiche2001DecomposableNN}.
What \textcite{Poon2011SumproductNA} termed \emph{Smoothness} was renamed \emph{smoothness} in subsequent work \textcite{Peharz2015OnTP} to avoid confusion with completeness in other sensesm and it requires that all children of every sum node share the same scope: for all $N, N' \in \mathrm{ch}(S)$, $\varphi(N)=\varphi(N')$ \cite{Poon2011SumproductNA, Peharz2015OnTP}.
As shown by \textcite{Peharz2015OnTP}, a PC that is both smooth and decomposable defines a valid (normalized) probability distribution, and computes the correct value of $P(\mathbf{X}=\mathbf{x})$ at the root for any assignment $\mathbf{x}$.

\emph{Sum-product networks} (SPNs), introduced by \textcite{Poon2011SumproductNA}, are the canonical instantiation of this framework: they are the first model family to demonstrate that deep, layered sum-product graphs can be learned directly from data while preserving tractable inference for the full distribution.
SPNs also subsume several older tractable model families, including thin junction trees, mixtures of trees, and certain classes of arithmetic circuits \cite{Poon2011SumproductNA, Darwiche2000ADA}.
The broader unifying term \emph{probabilistic circuit} was subsequently formalized by \cite{Choi2020ProbabilisticCA} to cover SPNs and related architectures such as PSDDs, cutset networks, and arithmetic circuits under a single framework;
we use the term PC throughout.

\subsection*{Tractable Probabilistic Models}\label{sec:pc-tractability}
% Marginalisation and conditioning are #P-hard in general models
% Show how decomposability + smoothness make marginalisation
%   a single downward pass — the property that §4.2 exploits
% Forward: CI testing requires P(V_i, V_j | Z); tractability is the prerequisite
The computational advantage of PCs rests on a fundamental contrast with general graphical models.
For arbitrary Bayesian networks over $d$ variables, \cite{Roth1993OnTH} showed that computing a marginal probability $P(V_i = v_i)$ is \#P-hard in the number of variables, so no polynomial-time algorithm is believed to exist under standard complexity-theoretic assumptions.
The hardness holds even for approximate marginalization under worst-case inputs \cite{Roth1993OnTH, Cooper1990TheCC}.
PCs escape this ceiling precisely because decomposability and smoothness turn marginalization into a structured pass through the circuit.

\paragraph{Marginalization.}
To marginalize a variable $X_j$ out of a PC, it suffices to replace the leaf distributions over $X_j$ with their normalized marginals (equivalently, to set the indicator for $X_j$ to the constant $1$) and evaluate the circuit upward from the leaves to the root \cite{Darwiche2000ADA, Poon2011SumproductNA}.
Decomposability ensures this substitution propagates correctly through every product node: since $\varphi(N) \cap \varphi(N') = \emptyset$ for the children of a product node, marginalizing $X_j$ affects only the child whose scope contains $X_j$, leaving all other children unchanged.
Smoothness ensures the substitution is consistent across all children of every sum node, since they share the same scope.
The result is that any marginal $P(\mathbf{X}_Q)$ over a query set $\mathbf{X}_Q \subseteq \mathbf{X}$ can be computed in time linear in the size of the circuit $|V|$, regardless of how many variables are summed out \cite{Darwiche2000ADA}.

\paragraph{Conditioning.}
A conditional distribution $P(\mathbf{X}_Q \mid \mathbf{X}_E = \mathbf{e})$ is obtained by evaluating the circuit with the indicator functions for $\mathbf{X}_E$ set to $\mathbf{e}$ (forward pass) and normalizing by the evidence probability $P(\mathbf{X}_E = \mathbf{e})$ obtained by marginalizing $\mathbf{X}_Q$ from the numerator circuit.
Both the forward pass and the normalization are linear-time operations under decomposability and smoothness \cite{Darwiche2000ADA, Choi2020ProbabilisticCA}.
The combination therefore supports full conditional inference in $\mathcal{O}(|V|)$ time per query.

\paragraph{Connection to CI Testing.}
As established in Section~\ref{sec:constraint-based}, the PC algorithm requires evaluating whether $V_i \perp\!\!\!\perp V_j \mid \mathbf{Z}$ for varying conditioning sets $\mathbf{Z}$.
A statistical test for this independence requires repeated evaluation of the joint conditional density $P(V_i, V_j \mid \mathbf{Z} = \mathbf{z})$ and the product of marginals $P(V_i \mid \mathbf{Z} = \mathbf{z}) \cdot P(V_j \mid \mathbf{Z} = \mathbf{z})$ across all samples.
Each of these quantities is a marginal or conditional query of the form above, and therefore runs in $\mathcal{O}(|V|)$ time in a smooth, decomposable PC.
This makes PCs a natural computational substrate for CI testing
in large-scale causal discovery pipelines: the otherwise intractable
density evaluations required by the PC algorithm become affordable
inner-loop operations. The existing federated CI test introduced
in Chapter~\ref{chap:fedcdh} bypasses this route entirely by
operating on kernel-based summary statistics rather than explicit
density queries \cite{Li2024FederatedCD}---a design choice that
avoids the tractability problem at the cost of expressive power.
Section~\ref{sec:federated-ci} addresses this gap by replacing
the kernel summary statistic with a PC-based density evaluator,
making the full joint conditional $P(V_i, V_j \mid \mathbf{Z})$
directly accessible as a tractable circuit query.

\subsection*{Scalable PC Architectures}\label{sec:pc-architectures}
% Einsum Networks: vectorised, GPU-friendly via tensor contraction
% RAT-SPNs: randomised scope partitioning for scalability
% FPC: fast inference for large-scale settings
% Forward: §4.3 uses FPC + RAT-SPN as leaf parameterisations

Standard structure-learning algorithms for PCs, such as LearnSPN
\cite{Gens2013LearningTS}, are accurate but scale poorly: they require centralized access to all data, and the recursive independence tests they perform to decide scope partitioning become the bottleneck for large $d$ or $n$.
Three architectural developments address these limitations and are directly relevant to the method developed in Chapter~\ref{chap:method}.

\paragraph{Einsum Networks (EiNets).}
\textcite{Peharz2020EinsumNF} observed that the forward pass of a PC can be re-expressed as a sequence of einsum operations over tensorized node arrays.
By grouping all sum and product nodes of a given depth into dense layers and executing the layer operations as batched tensor contractions, EiNets achieve GPU-friendly parallelism that avoids the irregular memory access patterns of
node-by-node evaluation.
This reduces the effective cost per inference pass by one to two orders of magnitude over unstructured node-by-node evaluation on comparable modern hardware \cite{Peharz2020EinsumNF}.
EiNets are one of the leaf parameterizations in the federated architecture introduced in Section~\ref{sec:pc-architecture-method}, alongside RAT-SPNs.

\paragraph{RAT-SPNs.}
A parallel scalability bottleneck is structure learning: greedy algorithms such as LearnSPN \cite{Gens2013LearningTS} search for the best variable-scope partition at each depth, requiring expensive CI tests on the training data.
\textcite{Peharz2019RandomSN} demonstrated that this search can be replaced by a random scope-partitioning strategy without sacrificing model quality.
A \emph{random-and-tensorized SPN} (RAT-SPN) is constructed by recursively and randomly dividing the root variable scope into two sub-scopes of equal size, repeating this process to a fixed depth $D$, and repeating the full tree-building procedure $R$ times (with independent random splits) to yield $R$ parallel scope-trees that are recombined at the root via a shared sum
node.
Each node array is tensorized and maps natively onto deep-learning frameworks, enabling GPU-parallel parameter estimation via gradient descent or expectation–maximization (EM) algorithm.
RAT-SPNs scale to millions of parameters while retaining the smoothness and decomposability guarantees required for tractable inference \cite{Peharz2019RandomSN}.
As Section~\ref{sec:pc-architecture-method} describes, RAT-SPNs are the structural template for the local PC models maintained at each client in the federated setting.

\subsection*{Distributional Conditions for Federated Settings}\label{sec:pc-federated-conditions}
% Formally define Mixture Marginals condition
% Formally define Cluster Independence condition
% Keep these as probabilistic definitions only — no FL mapping here
% The FL topology ↔ PC node duality is argued in §4.4, not here
% Cite: Seng et al. (2025)

Learning a globally valid PC from a distributed dataset requires two probabilistic conditions, formalized by \cite{Seng2025ScalingPC}, that ensure
the information held across clients is jointly sufficient to recover the global joint distribution $p(\mathbf{X})$.
Although Assumptions~\ref{ass:mixture-marginals} and~\ref{ass:cluster-independence} are stated in purely distributional terms, their motivation arises from settings in which a dataset $\mathbf{D}$ is partitioned across clients that cannot exchange raw observations.
Such settings are formalized as \emph{federated learning} (FL) in
Section~\ref{sec:fl}; the terminology is introduced here only to anchor the intuition behind each condition.
The connection to the FL topology --- specifically the duality between sum/product nodes and horizontal/vertical data partitioning --- is developed in Section~\ref{sec:pc-fl-duality}.

\begin{assumption}[Mixture Marginals {\cite[Assumption~1]{Seng2025ScalingPC}}]%
\label{ass:mixture-marginals}
Let $\mathbf{X} = \bigcup_{c \in C} \mathbf{X}_c$ be the union of all client variable sets, and let $\mathbf{X}_S \subseteq \mathbf{X}$ be any subset shared across clients with complement $\mathbf{X}_{S^-} = \mathbf{X} \setminus \mathbf{X}_S$.
There exists a joint distribution $p$ and a prior $q$ over a latent variable $L$ such that
\[
  \int_{\mathbf{X}_{S^-}} p(\mathbf{x})\,d\mathbf{x}
  = \sum_{l \in \mathcal{L}} q(L = l)\cdot p_S(\mathbf{x}_S \mid L = l)
  \quad \text{for all } \mathbf{x}_S \in \mathbf{X}_S,
\]
where each $p_S$ is defined over $\mathbf{X}_S$.
\end{assumption}

Intuitively, Assumption~\ref{ass:mixture-marginals} requires that the marginal of $p$ on any shared feature subset is representable as a finite mixture of the individual client distributions restricted to those features.
Without this condition, the clients' partitions would not collectively carry enough statistical information to reconstruct $p(\mathbf{X})$, rendering federated learning of the joint distribution ill-posed \cite{Seng2025ScalingPC}.
Note that this condition is the distributional counterpart of the smoothness requirement for sum nodes: sum nodes implement a convex combination over children with identical scopes (Eq.~\eqref{eq:sum-node}), and Mixture Marginals ensures this mixture has a valid global interpretation.

\begin{assumption}[Cluster Independence {\cite[Assumption~2]{Seng2025ScalingPC}}]%
\label{ass:cluster-independence}
Given disjoint sets of random variables $\mathbf{X}_1, \ldots, \mathbf{X}_n$ and a joint distribution $p(\mathbf{X}_1, \ldots, \mathbf{X}_n)$, there exists a latent variable $L$ with prior $q$ such that
\[
  p(\mathbf{X}_1, \ldots, \mathbf{X}_n)
  = \sum_{l \in \mathcal{L}} q(L = l)\prod_{i=1}^{n} p(\mathbf{X}_i \mid L = l).
\]
\end{assumption}

Assumption~\ref{ass:cluster-independence} extends conditional independence to the multi-client feature partition: given the latent cluster variable $L$, variables held by different clients are mutually independent.
This allows cross-client statistical dependencies to be captured through the shared latent $L$ — essentially a cluster selector — without requiring any direct exchange of raw features \cite{Seng2025ScalingPC}.
The assumption is strictly more expressive than a plain product distribution and is a standard modelling device in the PC literature, underpinning the validity of structure-learning algorithms such as LearnSPN that split scopes based on estimated independence \cite{Gens2013LearningTS}.
As shown by \textcite{Seng2025ScalingPC}, Assumption~\ref{ass:cluster-independence} is the distributional counterpart of the decomposability requirement for product nodes: product nodes implement a factorization over disjoint scopes (Eq.~\eqref{eq:product-node}), and Cluster Independence ensures this factorization — conditioned on the latent $L$ — accurately reflects the joint distribution when features reside on separate clients.
Together, Assumptions~\ref{ass:mixture-marginals}
and~\ref{ass:cluster-independence} are the minimal probabilistic prerequisites for the federated PC training procedures described in Chapter~\ref{chap:method}; their implications for the FL communication topology are analyzed in Section~\ref{sec:pc-fl-duality}.

% ─────────────────────────────────────────────────────────────────────────────
% §2.4 — Federated Learning
% Purpose of this section:
%   §3.1&2— reader understands the three FL partitioning regimes so the
%            sum/product ↔ horizontal/vertical mapping in FCs is legible
%   §4.4  — reader appreciates why unifying all three settings in one
%            framework is non-trivial, motivating FCs as a contribution
%   §4.6  — reader has the MPC / HE / DP vocabulary for the privacy analysis
% ─────────────────────────────────────────────────────────────────────────────
\section{Federated Learning}\label{sec:fl}

\subsection*{Setting and Motivation}
% One paragraph: why FL exists, the data sovereignty problem

Federated learning (FL) is a distributed machine learning paradigm in which a set of clients, aim to collaboratively learn a model without sharing their raw data \cite{McMahan2016CommunicationEfficientLO}. This setup directly addresses the data sovereignty problem: institutional and legal constraints routinely prohibit centralization of raw observations even when collaborative inference would be statistically beneficial. In the biomedical domain, neuroimaging records — such as the fMRI hippocampal time-series used in the FedCDH evaluation \cite{Poldrack2015LongtermNA} — fall under patient-privacy regulations (e.g., GDPR in the EU or HIPAA in the US) that restrict cross-institutional data transfer. In semiconductor and process manufacturing, equipment sensor logs encode proprietary process recipes and yield signals that constitute core intellectual property, making centralization contractually or competitively infeasible \cite{Vukovic2022CausalDI}. In both cases, inferring causal structure collaboratively across sites would yield more statistically reliable graphs than any single institution's sample size could support alone, yet the data cannot leave its origin.


\subsection*{Data Partitioning Topologies}
How the dataset $\mathbf{D} \in \mathbb{R}^{n \times d}$ is partitioned across clients determines which learning algorithms are applicable and what can be communicated without violating data sovereignty.
Three regimes are distinguished depending on whether the partition is along the sample axis, the feature axis, or both \cite{Seng2025ScalingPC}; see Table~\ref{tab:fl-topologies} for a comparison.
The following definitions invoke the \emph{Mixture Marginals}
(Assumption~\ref{ass:mixture-marginals}) and \emph{Cluster Independence} (Assumption~\ref{ass:cluster-independence}) conditions established in Section~\ref{sec:pc-federated-conditions}, which together ensure the distributed data jointly suffices to recover $p(\mathbf{X})$.

% TODO: Maybe a Schematic Figure showing the three partition modes side by side (samples on rows, features on columns, clients as coloured blocks) would be worth approximately half a page and would make §4.4's later argument about sum/product nodes much easier to follow.

\begin{definition}[Horizontal FL {\cite[Def.~1]{Seng2025ScalingPC}}]
Assume a set of samples $\mathbf{D}_c \sim p_c$ on each client $c \in C$, a joint distribution $p$ satisfying the Mixture Marginals condition, and that $\mathbf{X}_c = \mathbf{X}_{c'}$ for all $c, c' \in C$ s.t.\ $c \neq c'$. Horizontal FL is defined as fitting a mixture $\hat{p} = \sum_{c \in C} q(c) \cdot \hat{p}_c$ such that $d(\hat{p}, p)$ and $d(p_c, \hat{p}_c)$ are minimal for all $c \in C$, where $d$ is a distance metric and $\hat{p}_c$ are local distribution estimates.
\end{definition}

\begin{definition}[Vertical FL {\cite[Def.~2]{Seng2025ScalingPC}}]
Assume a set of samples $\mathbf{D}_c \sim p_c$ on each client $c \in C$, a joint distribution $p$ satisfying both the Mixture Marginals and Cluster Independence conditions, and that $\mathbf{X}_c \cap \mathbf{X}_{c'} = \emptyset$ for all $c, c' \in C$ s.t.\ $c \neq c'$. Vertical FL is defined as estimating a joint distribution $\hat{p}$ such that $d(p, \hat{p})$ is minimal and $\sum_{X \setminus X_c} \hat{p}(x) = \hat{p}_c(x)$ for all $x \in X$, where $d$ is a distance metric, $\hat{p}_c$ are estimates of client distributions, and $\sum$ denotes integration for continuous and summation for discrete variables.
\end{definition}

Hybrid FL is defined as a combination of Definitions 1 and 2, permitting partial overlap in both sample and feature dimensions \cite{Wen2022ASO}. Prior to unified probabilistic frameworks, reconciling all three settings within a single method was considered non-trivial in the FL community \cite{ChaveroDiez2026FederatedLF, Seng2025ScalingPC}.

\subsection*{Standard Training Protocols}
% FedAvg for horizontal; SplitNN for vertical; communication costs
% One table or diagram showing the three topologies

In horizontal FL, the dominant training protocol is FedAvg \cite{McMahan2016CommunicationEfficientLO}, in which clients iteratively optimize local objectives and send model parameters to a central server for aggregation. For a model with $M$ parameters, $|C|$ clients, and $K_r$ communication rounds, this incurs $\mathcal{O}(M \cdot |C| \cdot K_r)$ network messages. However, model averaging requires all clients to share the same model structure, which is only guaranteed when feature spaces are aligned — a condition that holds in horizontal FL but not in vertical FL \cite{Seng2025ScalingPC}.

In vertical FL, SplitNN-like architectures are the standard choice where clients share intermediate feature representations rather than model parameters \cite{Ceballos2020SplitNNdrivenVP}; this incurs $\mathcal{O}(E \cdot |C| \cdot F \cdot n)$ messages for $E$ training epochs, with each client outputting a feature representation of size $F$ over $n$ samples \cite{Seng2025ScalingPC}.

In hybrid FL, communication cost combines both contributions:
$\mathcal{O}(|R_s| \cdot C_h + |R_v| \cdot C_v)$, where $R_s$ denotes the set of feature subspaces shared across multiple clients (handled horizontally with per-subspace cost $C_h$), $R_v$ the feature subspaces residing on a single client (handled
vertically with per-subspace cost $C_v$), and the two terms are summed over all disjoint subspaces \cite{Seng2025ScalingPC}.

\begin{table}
\centering
\caption{The three FL data partitioning topologies.}
\label{tab:fl-topologies}
\begin{tabular}{lccc}
\toprule
 & \textbf{Horizontal FL} & \textbf{Vertical FL} & \textbf{Hybrid FL} \\
\midrule
Same samples across clients & \ding{55} & \ding{51} & partial \\
Same features across clients & \ding{51} & \ding{55} & partial \\
Standard protocol & FedAvg & SplitNN & -- \\
Comm. cost (dominant term) & $\mathcal{O}(M\cdot|C|\cdot K_r)$ & $\mathcal{O}(E\cdot|C|\cdot F \cdot n)$ & $\mathcal{O}(|R_s|\cdot C_h + |R_v| \cdot C_v)$ \\
\bottomrule
\end{tabular}
\end{table}

\subsection*{Privacy Mechanisms}
% Secure MPC, homomorphic encryption, differential privacy — 2-3 sentences each
% Keep brief: full privacy analysis is in §4.6
The data partitioning constraint in FL provides a first layer of protection, but sharing even summary statistics or aggregated local quantities does not fully eliminate the risk of information exposure \cite{Li2024FederatedCD}. FL pipelines are therefore commonly augmented with cryptographic or statistical mechanisms that provide stronger, quantifiable guarantees \cite{Dwork2006CalibratingNT, Dwork2014TheAF}.

\textbf{Secure multiparty computation (MPC)} allows a set of parties to jointly compute a function of their inputs while keeping those inputs private from all other parties \cite{Cramer2015}. In FL, this enables clients to contribute to a global aggregate on the server without any individual client's input being observable in isolation \cite{Li2024FederatedCD}.

\textbf{Homomorphic encryption (HE)} enables complex mathematical operations to be performed on encrypted data without requiring decryption \cite{Doan2023ASO}. A client can therefore transmit an encrypted local statistic that the server aggregates entirely in the ciphertext domain, though at significant computational cost \cite{Li2024FederatedCD}.

\textbf{Differential privacy (DP)} provides a formal statistical guarantee that the output of a computation does not reveal whether any particular record was present in the input \cite{Dwork2014TheAF}. Unlike MPC and HE, which protect data in transit or during computation, DP operates by injecting calibrated noise into outputs, bounding the privacy loss that an adversary can infer \cite{Dwork2006CalibratingNT}. It is parameterized by $(\varepsilon, \delta)$-DP, where $\varepsilon$ controls the privacy-loss budget and $\delta$ bounds the probability of that budget being exceeded; smaller $\varepsilon$ implies stronger privacy at the cost of greater noise injection \cite{Dwork2014TheAF}.

A full analysis of the privacy-utility trade-offs under these mechanisms in the context of this thesis is deferred to Section~\ref{sec:privacy}.


\section{Federated Causal Discovery from Heterogeneous Data (FedCDH)}\label{sec:FedCDH}
\subsubsection**{Heterogeneous Data and the Surrogate Variable}

\subsection*{Kernel-based Conditional Independence Testing and the Hilbert-Schmidt Independence Criterion }

\subsection*{Formal Setup and Assumptions}

\subsection*{The Federated Conditional Independence Test (FCIT)}\label{sec:pc-based-ci}

\subsection*{The Federated Independent Change Principle (FICP)}


\subsection*{Summary Statistics Protocol and Sufficiency }


\subsection*{Existing Federated Causal Discovery Methods and Open Gaps}\label{sec:federated-baselines}



\section{Federated Probabilistic Circuits (FedPCs).}\label{subsec:fedPCs}
\textcite{Seng2025ScalingPC} introduced \emph{federated circuits} (FCs) and their probabilistic instantiation \emph{federated PCs} (FedPCs) as a framework for learning PCs when data are distributed across \emph{clients} --- independent data owners that cannot share raw
observations.
A FedPC is defined as a federated circuit in which each leaf node is a local density estimator (itself a PC, e.g.\ an EiNet or RAT-SPN) and each internal node is either a sum or a product node executed as a distributed computation across the participating machines \cite{Seng2025ScalingPC}.
Crucially, FedPCs preserve the structural validity properties of standard PCs: the combined circuit is smooth and decomposable by construction, so all tractability guarantees for marginalization and conditioning carry over to the federated setting.
Training proceeds via Algorithm 1 of \textcite{Seng2025ScalingPC} (one-pass algorithm) in which clients learn their local leaf models independently and contribute only summary statistics or scope identifiers to the server, keeping raw data private \cite{Seng2025ScalingPC}.
Empirically, FedPCs match the log-likelihood of centralized PCs on tabular benchmarks while achieving training speedups of $4$--$14\times$ through parallelism \cite{Seng2025ScalingPC}.
The FedPC framework is the direct technical antecedent of the method proposed in Chapter~\ref{chap:method}, where it is adapted to the causal discovery objective.


% \chapter{Federated Causal Discovery from Heterogeneous Data (FedCDH)}\label{chap:foundation}
% \subsubsection**{Heterogeneous Data and the Surrogate Variable}

\subsection*{Kernel-based Conditional Independence Testing and the Hilbert-Schmidt Independence Criterion }

\subsection*{Formal Setup and Assumptions}

\subsection*{The Federated Conditional Independence Test (FCIT)}\label{sec:pc-based-ci}

\subsection*{The Federated Independent Change Principle (FICP)}


\subsection*{Summary Statistics Protocol and Sufficiency }


\subsection*{Existing Federated Causal Discovery Methods and Open Gaps}\label{sec:federated-baselines}


\chapter{Proposed Method: Probabilistic Circuit Federated Causal Discovery (PC-FedCD)}\label{chap:method}
\section{Method Overview and Architecture}


\section{PC-Based Conditional Independence Testing}\label{sec:federated-ci}


\section{Local Structure Estimation via FPC and RAT-SPNs}\label{sec:pc-architecture-method}


\section{Federated Circuits Integration: Horizontal, Vertical, and Hybrid Data Splits}\label{sec:pc-fl-duality}


\section{One-Pass Training for Federated Causal Discovery}


\section{Privacy Analysis and Communication Efficiency}\label{sec:privacy}


\chapter{Experimental Evaluation}\label{chap:experiment}
\section{Experimental Setup}
\lipsum[2]

\section{Causal Accuracy Against Baselines}
\lipsum[2]

\section{Scalability Analysis}
\lipsum[2]

\section{Horizontal, Vertical, and Hybrid Federated Settings}
\lipsum[2]

\section{Privacy-Utility Trade-off}
\lipsum[2]

\section{Ablation Study}
\lipsum[2]

\section{Real-World Applications}
\lipsum[2]


\chapter{Discussion}\label{chap:discussion}
\section{Strengths over FedCDH and Other Baselines}
\lipsum[2]

\section{Limitations and Failure Cases}
\lipsum[2]

\section{Privacy Guarantees and Communication Overhead in Practice}
\lipsum[2]

\section{Future Directions}
\lipsum[2]


\chapter{Conclusion}\label{chap:conclusion}
\input{chapters/7_conclusion}

\backmatter
\printbibliography[heading=bibintoc, title={Bibliography}]

\appendix
\chapter{Appendix A}
\lipsum[4]

\appendix
\chapter{Appendix B}
\lipsum[4]

\end{document}
